% Lösungen Einheit 1 von 3 – Irrationale Zahlen und Zahlbereiche (zusammenbau v0.1)
\einheitenkopf{Einheit 1 von 3 · Irrationale Zahlen und Zahlbereiche}

%% TODO Lösung der Erklärzeile 10g) fehlt in der Bank
\erg{10}{a) geht auf \quad b) $0{,}625$ – abbrechend \quad c) $0{,}\overline{4}$ ($\approx 0{,}444$) – periodisch \quad d) $0{,}275$ – abbrechend \quad e) $0{,}8\overline{3}$ ($\approx 0{,}833$) – periodisch \quad f) $0{,}\overline{18}$ ($\approx 0{,}182$) – periodisch \quad g) \ldots \quad h) rational, denn $\sqrt{169} = 13$ \quad i) $\sqrt{5} \approx 2{,}236$ \quad j) ℤ, ℚ und ℝ \quad k) $2{,}4$ und $2{,}5$ ($\sqrt{6} \approx 2{,}45$) \quad l) $2{,}8$; $\sqrt{8}$; $2\frac{5}{6}$ (denn $\sqrt{8} \approx 2{,}828$ und $2\frac{5}{6} \approx 2{,}833$) \quad m) $a = \sqrt{7}\,\text{cm} \approx 2{,}65\,\text{cm}$ \quad n) $7$ – rational \quad o) $1{,}96$ und $2{,}25$, also $1{,}4 < \sqrt{2} < 1{,}5$ \quad p) irrational, weil der Radikand keine Quadratzahl ist; $\sqrt{40} \approx 6{,}32$; nur ℝ}

10k)

\zahlenstrahl[xmin=2,xmax=3,xstep=0.1,karo=1]{2.449/\sqrt{6}}

\erg{11}{a) $\frac{9}{20}$}
\erg{12}{a) Ja: Jede ganze Zahl lässt sich als Bruch mit Nenner eins schreiben.}
\erg{13}{a) Die Anzeige ist nur ein Näherungswert, $\sqrt{3}$ ist irrational; richtig: $\sqrt{3} \approx 1{,}732$. \quad b) Wäre $\sqrt{2}$ ein gekürzter Bruch, müsste das Quadrat des Zählers doppelt so groß sein wie das des Nenners; dann wären Zähler und Nenner beide gerade – der Bruch wäre nicht gekürzt, ein Widerspruch. \quad c) $\sqrt{8}$ \quad d) $a = \sqrt{6}\,\text{m} \approx 2{,}45\,\text{m}$}
